% GENERATED FILE. Edit canonical lessons, then run npm run generate:books.
% canonical-source-hash: fnv1a64:a14908fb9b0cf048
% canonical-lessons: HI-C234-baya, HI-C234-sidha, HI-C234-ulta, HI-C234-bic, HI-C234-kona

\chapter{Left, Straight, Upside down, The middle, A corner}
\label{ch:hi-left-straight-upside-down-the-middle-a-corner}
\hlchaptermodality{234}

\begin{chapteropening}
\textbf{By the end of this chapter:} \emph{I can say left, straight, upside down, the middle and a corner.}

Complete the last lesson of chapter 234: i can say left, straight, upside down, the middle and a corner.
\end{chapteropening}

\section[\texorpdfstring{\hi{बायाँ} (bāyā̃)}{बायाँ (bāyā̃)}]{\hi{बायाँ} (bāyā̃) — left (side)}

\label{lesson:HI-C234-baya}



\begin{quote}
Before the new one: say the Hindi for west, then the Hindi for right (side).
\end{quote}

\subsection*{You'll want to know: \hi{बायाँ}}
\textbf{\hi{बायाँ}} — \emph{bāyā̃} — \textquotedblleft{}left (side)\textquotedblright{}.

\begin{grammarlens}[title={What you've built}]
One more word for this chapter.
\end{grammarlens}

\subsection*{Your turn}
\begin{itemize}
\raggedright
  \item \emph{Say it:} \emph{bāyā̃}
  \item \emph{Say it:} \emph{bāyā̃}, once more
  \item \emph{Say it:} say \emph{bāyā̃}
  \item \emph{Recall:} say the Hindi for west, then the Hindi for right (side), then say \emph{bāyā̃} again
\end{itemize}

\begin{tcolorbox}[breakable,colback=teal!4,colframe=teal!35!black,title={Before you move on}]
What does \hi{बायाँ} mean? (\textquotedblleft{}Left (side)\textquotedblright{}.) Say it once more.
\end{tcolorbox}
\section[\texorpdfstring{\hi{सीधा} (sīdhā)}{सीधा (sīdhā)}]{\hi{सीधा} (sīdhā) — straight, direct}

\label{lesson:HI-C234-sidha}



\begin{quote}
Before the new one: say the Hindi for right (side), then the Hindi for left.
\end{quote}

\subsection*{You'll want to know: \hi{सीधा}}
\textbf{\hi{सीधा}} — \emph{sīdhā} — \textquotedblleft{}straight, direct\textquotedblright{}.

\begin{grammarlens}[title={What you've built}]
One more word for this chapter.
\end{grammarlens}

\subsection*{Your turn}
\begin{itemize}
\raggedright
  \item \emph{Say it:} \emph{sīdhā}
  \item \emph{Say it:} \emph{sīdhā}, once more
  \item \emph{Say it:} say \emph{sīdhā}
  \item \emph{Recall:} say the Hindi for right (side), then the Hindi for left, then say \emph{sīdhā} again
\end{itemize}

\begin{tcolorbox}[breakable,colback=teal!4,colframe=teal!35!black,title={Before you move on}]
What does \hi{सीधा} mean? (\textquotedblleft{}Straight, direct\textquotedblright{}.) Say it once more.
\end{tcolorbox}
\section[\texorpdfstring{\hi{उल्टा} (ulṭā)}{उल्टा (ulṭā)}]{\hi{उल्टा} (ulṭā) — upside down, reversed}

\label{lesson:HI-C234-ulta}



\begin{quote}
Before the new one: say the Hindi for left, then the Hindi for straight.
\end{quote}

\subsection*{You'll want to know: \hi{उल्टा}}
\textbf{\hi{उल्टा}} — \emph{ulṭā} — \textquotedblleft{}upside down, reversed\textquotedblright{}.

\begin{grammarlens}[title={What you've built}]
One more word for this chapter.
\end{grammarlens}

\subsection*{Your turn}
\begin{itemize}
\raggedright
  \item \emph{Say it:} \emph{ulṭā}
  \item \emph{Say it:} \emph{ulṭā}, once more
  \item \emph{Say it:} say \emph{ulṭā}
  \item \emph{Recall:} say the Hindi for left, then the Hindi for straight, then say \emph{ulṭā} again
\end{itemize}

\begin{tcolorbox}[breakable,colback=teal!4,colframe=teal!35!black,title={Before you move on}]
What does \hi{उल्टा} mean? (\textquotedblleft{}Upside down, reversed\textquotedblright{}.) Say it once more.
\end{tcolorbox}
\section[\texorpdfstring{\hi{बीच} (bīc)}{बीच (bīc)}]{\hi{बीच} (bīc) — the middle}

\label{lesson:HI-C234-bic}



\begin{quote}
Before the new one: say the Hindi for straight, then the Hindi for upside down.
\end{quote}

\subsection*{You'll want to know: \hi{बीच}}
\textbf{\hi{बीच}} — \emph{bīc} — \textquotedblleft{}the middle\textquotedblright{}.

\begin{grammarlens}[title={What you've built}]
One more word for this chapter.
\end{grammarlens}

\subsection*{Your turn}
\begin{itemize}
\raggedright
  \item \emph{Say it:} \emph{bīc}
  \item \emph{Say it:} \emph{bīc}, once more
  \item \emph{Say it:} say \emph{bīc}
  \item \emph{Recall:} say the Hindi for straight, then the Hindi for upside down, then say \emph{bīc} again
\end{itemize}

\begin{tcolorbox}[breakable,colback=teal!4,colframe=teal!35!black,title={Before you move on}]
What does \hi{बीच} mean? (\textquotedblleft{}The middle\textquotedblright{}.) Say it once more.
\end{tcolorbox}
\section[\texorpdfstring{\hi{कोना} (konā)}{कोना (konā)}]{\hi{कोना} (konā) — a corner}

\label{lesson:HI-C234-kona}



\begin{quote}
Before the new one: say the Hindi for upside down, then the Hindi for the middle.
\end{quote}

\subsection*{You'll want to know: \hi{कोना}}
\textbf{\hi{कोना}} — \emph{konā} — \textquotedblleft{}a corner\textquotedblright{}.

\begin{grammarlens}[title={What you've built}]
One more word for this chapter.
\end{grammarlens}

\subsection*{Your turn}
\begin{itemize}
\raggedright
  \item \emph{Say it:} \emph{konā}
  \item \emph{Say it:} \emph{konā}, once more
  \item \emph{Say it:} say \emph{konā}
  \item \emph{Recall:} say the Hindi for upside down, then the Hindi for the middle, then say \emph{konā} again
\end{itemize}

\begin{tcolorbox}[breakable,colback=teal!4,colframe=teal!35!black,title={Before you move on}]
What does \hi{कोना} mean? (\textquotedblleft{}A corner\textquotedblright{}.) Say it once more.
\end{tcolorbox}
